\documentclass{article}

\begin{document}
    This is the most bare bones possible set up for a \LaTeX\footnote{Pronounced LAY-tech or LAH-tech, as it is based off of TeX typesetting system and the `TeX' was originally $\tau\epsilon\chi$, the foundation for the English word `technique'.} document possible. This is a file dedicated to the pure fundamentals of what LaTeX is, how to begin a simple document, how to typescript some mathematics, and a bit of standard formatting philosophy sprinkled in.
    \begin{itemize}
        \item \verb|\documentclass| and \verb|\begin{document}|$\dots$\verb|\end{document}|
        \item 
        \item Typesetting basic math
    \end{itemize}
    \section{The Core}
    The core of a LaTeX file is the: \verb|\documentclass{|$\bullet$\verb|}|
    \section{Typesetting math}
    Paragraphs should be separated by white lines, this is a standard practice and makes your document more legible before compilation. We will cover other ways of controlling white space later. For writing things in a mathematical style, you need to enter `math mode', which is controlled by surrounding things in a pair of \$'s. Within math mode, you unlock commands that would otherwise throw an error in LaTeX. For example: $x = \ln $
\end{document}
